\documentclass[output=paper]{langscibook}
\ChapterDOI{10.5281/zenodo.22880306}
\title{\textit{Hovuhovu iobo} – Digging out sea-sausages}
\author{Helen Kobaa Magum and Shalom Magum}
\abstract{}
\begin{document}
\maketitle\label{chap:39}
\begin{figure}[H]
    \includegraphics[trim=0 0 0 0,clip,width=\linewidth]{figures/39_1.jpg}
    \caption{\textit{Hovuhovu iobo} -- Digging out sea-sausages}
    \label{fig:39_1}
\end{figure}
\begin{xltabular}{\textwidth}{lXX}
	(\ref{ex:text32-1})&\textit{\textbf{Hovuhovu iobo}}&\textbf{Digging out sea-sausages}\\
	(\ref{ex:text32-2})&\textit{O iobo a peha taba vaa tahii to aniani raara a ta Teapu.}&The sea-sausage is an animal of the sea that we Teop people eat.\\
	(\ref{ex:text32-3})&\textit{A taba bona na teitei batana paano kasuana ei komana tahii tea maa hum kakaamuru.}&This animal stays underground in the sea in places with white sand.\\
	(\ref{ex:text32-4})&\textit{Beara tau nao tea vagana iobo, eara toro dee o pusu hiovu, a sinaka ae a bonaa, eara re nao tea masi.}&Before we go to catch sea-sausages, we must take a digging stick, a \textit{sinaka} stick and a fishing basket, then we go to the reef.\\
	(\ref{ex:text32-5})&\textit{A bonaa he na vabibiinoo ni raara raparaara.}&The fishing basket we carry on our side.\\
\end{xltabular}
\begin{figure}[H]
    \includegraphics[trim=0 0 0 0,clip,width=.8\linewidth]{figures/39_2.png}
    \caption{\textit{Pusu hiovu} -- Digging stick}
    \label{fig:39_2}
\end{figure}
\begin{figure}[H]
    \includegraphics[trim=0 0 0 0,clip,width=.8\linewidth]{figures/39_3.png}
    \caption{\textit{Sinaka}}
    \label{fig:39_3}
\end{figure}
\begin{figure}[H]
    \captionsetup[subfigure]{oneside,margin={0cm,0cm}}
    \subfloat[]{\includegraphics[trim=0 0 0 0,clip,width=.44\linewidth]{figures/39_4.png}\label{fig:39_4}}\hfil
    \subfloat[]{\includegraphics[trim=0 150 0 150,clip,width=.54\linewidth]{figures/39_5.jpg}\label{fig:39_5}}
    \caption{(a) \textit{Bonaa vabibiinoo} -- Fishing basket, (b) \textit{Iobo} -- Sea-sausage}
\end{figure}
\begin{xltabular}{\textwidth}{lXX}
	(\ref{ex:text32-6})&\textit{Beara suguna teo hum hovuhovu iobo, eara repaa taneo tea tara rakerake bata riori.}&When we arrive at the place for digging out sea-sausages, we start searching for them.\\
	(\ref{ex:text32-7})&\textit{Beara tara a sii puhuku kakaamuru, eara repaa kona voen o iobo na tei nana.}&When we see a little heap of white sand, we know that there is a sea-sausage.\\
	(\ref{ex:text32-8})&\textit{Eara repaa hovu hopo ni nao a ta hiovu, beara hogo bata ni naora ta meha papa kuriraara.}&Then we stick the digging-stick inside (the heap), while we are pushing our other hand inside.\\
	(\ref{ex:text32-9})&\textit{Eove to kahi ato naona bono iobo bea ta hiovu gae hiava maa bona.}&It will get hold of the sea-sausage when the digging stick lifts it up.\\
	(\ref{ex:text32-10})&\textit{Beara mei rae tea vaobete e komana bonaa, eara toro pisin roho e, eara repaa vaobete mana e komana bonaa.}&Before we put it into the fishing basket, we must turn the sea-sausage inside out, then we eventually put it into the basket.\\
	(\ref{ex:text32-11})&\textit{Eara pasi paku bata koa vo nao nom teebona, ore antee komana beo manin vagana tean kavara ge ei beo puhana puhana maa.}&We will do it like this, until you don't feel like fishing anymore or the high tide has come in.\\
    (\ref{ex:text32-12})&\textit{A kiu vaa tea tatasu iobo na pakupaku vaha raara kasuana en vaan.}&The work of beating the sea-sausages, we do it again on the beach here in the village.\\
	(\ref{ex:text32-13})&\textit{Eara na tasu vatamee ae vamuraka ni rara o iobo, tea tabae o iobo o taba o torisi.}&We beat the sea-sausage flexible and soft, because the sea-sausage is a tough thing.\\
	(\ref{ex:text32-14})&\textit{Tea tatasu he eara na ibun rara amaa iobo tea sinaka, eara repaa tatasu riori bona ta hiovu.}&For the beating we put some sea-sausages on the \textit{sinaka} stick, and then we beat them with the digging stick.\\
	(\ref{ex:text32-15})&\textit{Beo iobo bebeahu ae taveraha, eara repaa kona voen paa muraka vai vuru, eara repaa antee tea tasu.}&When the sea-sausage is long and flattened, we know that it is soft now so that we can stop beating it.\\
\end{xltabular}
\begin{figure}[H]
    \captionsetup[subfigure]{oneside,margin={0cm,0cm}}
    \subfloat[]{\includegraphics[trim=0 0 0 0,clip,width=.49\linewidth]{figures/39_6.jpg}\label{fig:39_6}}\hfil
    \subfloat[]{\includegraphics[trim=0 0 0 0,clip,width=.49\linewidth]{figures/39_7.jpg}\label{fig:39_7}}
    \caption{(a) \textit{Ibun iobo} -- Stick with sea-sausages, (b) \textit{Tatasu iobo} -- Beating sea-sausages}
\end{figure}
\begin{figure}[H]
    \captionsetup[subfigure]{oneside,margin={0cm,0cm}}
    \subfloat[]{\includegraphics[trim=0 0 0 0,clip,width=.49\linewidth]{figures/39_8.jpg}\label{fig:39_8}}\hfil
    \subfloat[]{\includegraphics[trim=0 0 0 0,clip,width=.49\linewidth]{figures/39_9.jpg}\label{fig:39_9}}
    \caption{(a) \textit{Geregerete} -- Tearing off the upper part, (b) \textit{Ahuhu} -- Squeezing out the water}
\end{figure}
\begin{xltabular}{\textwidth}{lXX}
	(\ref{ex:text32-16})&\textit{Bea tatasu kavara, eara repaa geregerete riori ae ahuhu.}&When the beating is finished, we tear off the upper parts and squeeze the water out (of the sea-sausage).\\
	(\ref{ex:text32-17})&\textit{Be kavara eara repaa voosu vaan.}&When it is finished, we return home.\\
	(\ref{ex:text32-18})&\textit{Tea nahu nae he na paku vonaen:}&As for their cooking, it is done as follows:\\
	(\ref{ex:text32-19})&\textit{Eara na vaobete naora amaa iobo komana nahu.}&We put the sea-saugages into the pot.\\
	(\ref{ex:text32-20})&\textit{Eara repaa koma hopo nia ruene, tasu a ta hiihii ae o hiiruu.}&Then we pour water into it, pound a piece of ginger and basil leaves.\\
	(\ref{ex:text32-21})&\textit{Eara vaobete hopo e komana nahu eara repaa vaagonono.}&We put them into the pot and then we let the water boil.\\
	(\ref{ex:text32-22})&\textit{Beve mahaka, eara kou an vai e, a huun nae he na hiroo bata raara.}&When they are done, we eat them, while drinking their liquid.\\
	(\ref{ex:text32-23})&\textit{Evehee toro vaorete bata mi raara bona kaukau, muu, kavakava gea tapeako.}&But we must eat them together with sweet potatoes, taro, cooked bananas or manioc.\\
\end{xltabular}
\end{document}